\section{Introduction}
\label{sec:introduction}

\PARstart{E}{motional} voice conversion (EVC) seeks to transform the perceived emotion of an utterance while preserving linguistic content and speaker identity.
This dual objective creates a central evaluation tension: a system can appear successful because target-emotion metrics improve, while speaker identity quietly drifts.
The standard EVC goal is therefore not only to make converted speech sound emotional, but to do so with minimum degradation of speaker characteristics. \cite{yang2022overview,zhou2022evctheoryesd}.

Most EVC evaluations report emotion-side and speaker-side measurements as separate aggregate scores.
That practice is useful, but it does not reveal how the two objectives interact when target-emotion strength is actively varied.
In particular, a single high-intensity operating point cannot show whether a method preserves speaker identity across weaker and stronger emotional settings.
This paper studies the hypothesis that emotional intensity exposes an underdiscussed speaker--emotion frontier: as target emotion becomes stronger, speaker preservation may remain stable, degrade smoothly, or collapse depending on the model, dataset, and conversion condition.

Instead of treating emotion accuracy and speaker similarity as independent numbers, we treat each intensity setting as an operating point.
The curve traced by these operating points is the object of analysis.
This framing makes the evaluation concrete: a method is characterized by the operating region it traces, the speaker-preservation cost of increasing emotion, and the dataset or conversion condition in which that frontier is observed.


We frame this primarily as a measurement problem, with a model-side mechanism for tracing controllable operating points.
The measurement contribution is an intensity-conditioned evaluation protocol that traces emotion-side metrics against speaker-side metrics across an intensity sweep.
The modeling contribution is a native-intensity StarGANv2-VC-based architecture designed to represent target-emotion strength explicitly, not a claim that speaker--emotion disentanglement is solved.
We additionally study speaker-related ablations, including speaker-contrastive and speaker-classifier variants, to test how speaker supervision reshapes the frontier.
We compare this model family against two baselines: StarGANv2-VC with inferred intensity and a CycleGAN spectrum-conversion baseline.

The main contributions of this work are:
\begin{itemize}
    \item An intensity-conditioned evaluation protocol for measuring the speaker--emotion frontier in EVC.
    \item A native-intensity StarGANv2-VC-based architecture for tracing controllable emotional conversion operating points.
    \item A comparison between native intensity, inferred-intensity StarGANv2-VC, and CycleGAN across ESD \cite{zhou2021seen} and CREMA-D \cite{cao2014cremad} datasets.
\end{itemize}

The remainder of this paper is organized as follows:
Section~\ref{sec:related-work} reviews EVC, emotion-intensity control, speaker preservation, and the missing frontier-style evaluation.
Section~\ref{sec:problem-formulation} defines inferred and native intensity sweeps and the proposed frontier measurement.
Section~\ref{sec:proposed-model} describes the native-intensity model and speaker-related ablations.
Section~\ref{sec:experimental-design} gives the datasets, baselines, metrics, and ablations.
Sections~\ref{sec:results}--\ref{sec:limitations} present results, discussion, and limitations.
